\subsection{Cross-Dataset Configuration Analysis}
\label{sec:cross-dataset-entropy}

The factor ablation (Section~\ref{sec:factor-ablation}) and weight grid search (Section~\ref{sec:weight-learning}) identify the two-factor configuration $(\alpha{=}0, \beta{=}0.5, \gamma{=}0.5, \delta{=}0)$ as optimal on ImageNet-100. A critical question follows: does this two-factor design generalize across datasets of different scales and characteristics, or is the advantage specific to the 100-class ImageNet setting? To answer this, we evaluate the four-factor \cags{} ($\alpha{=}0.3, \beta{=}0.3, \gamma{=}0.2, \delta{=}0.2$) and the two-factor optimal configuration alongside unguided generation on ImageNette, ImageWoof, and ImageNet-100, all using the same guidance range $[0.0, 0.06]$, sigmoid parameters $(s{=}3.0, c_0{=}0.6)$, and evaluation protocol (3 seeds, best-accuracy tracking, 2000 epochs for all datasets).

\Cref{tab:cross-dataset-config} presents the results. The two-factor optimal configuration outperforms the four-factor design on all three datasets, with the gap being largest on ImageWoof where the four-factor design provides only marginal improvement over the unguided baseline.

\begin{table}[t]
\centering
\caption{Cross-dataset comparison of unguided, four-factor \cags{}, and the two-factor optimal configuration. All configurations use the same guidance range $[0.0, 0.06]$, sigmoid parameters $(3.0, 0.6)$, and evaluation protocol (3 seeds, best-accuracy tracking, 2000 epochs). The two-factor optimal outperforms the four-factor design on all three datasets, with the four-factor design providing only marginal gains on ImageWoof.}
\label{tab:cross-dataset-config}
\begin{tabular}{llccc}
\toprule
\textbf{Dataset} & \textbf{Method} & \textbf{Top-1 (\%)} & \textbf{$\Delta$ vs.\ unguided} & \textbf{$\Delta$ opt.\ vs.\ 4-fac.} \\
\midrule
\multirow{3}{*}{ImageNette (10-cls)} & Unguided & $49.99 \pm 0.22$ & --- & --- \\
 & \cags{} (4-factor) & $52.64 \pm 0.28$ & $+2.65$ & --- \\
 & \cags{} (2-factor opt.) & $\mathbf{55.27 \pm 0.68}$ & $\mathbf{+5.28}$ & $+2.63$ \\
\midrule
\multirow{3}{*}{ImageWoof (10-cls)} & Unguided & $31.75 \pm 0.06$ & --- & --- \\
 & \cags{} (4-factor) & $31.99 \pm 0.46$ & $+0.24$ & --- \\
 & \cags{} (2-factor opt.) & $\mathbf{32.87 \pm 0.29}$ & $\mathbf{+1.12}$ & $+0.88$ \\
\midrule
\multirow{3}{*}{ImageNet-100 (100-cls)} & Unguided & $22.79 \pm 0.18$ & --- & --- \\
 & \cags{} (4-factor) & $25.29 \pm 0.27$ & $+2.50$ & --- \\
 & \cags{} (2-factor opt.) & $\mathbf{25.70 \pm 0.07}$ & $\mathbf{+2.91}$ & $+0.41$ \\
\bottomrule
\end{tabular}
\end{table}

\paragraph{The four-factor design is ineffective on fine-grained datasets.}
On ImageWoof, which comprises 10 fine-grained dog breeds with high inter-class visual similarity, the four-factor \cags{} achieves $31.99 \pm 0.46\%$---only $+0.24\%$ above the unguided baseline ($31.75 \pm 0.06\%$, $+0.8\%$ relative), a negligible improvement. This stagnation occurs because the separability factor ($\delta{=}0.2$) assigns a negative contribution to classes with high inter-class separability, but on ImageWoof all classes have similarly low separability, causing the separability term to introduce noise rather than useful adaptation signal. The mode count factor ($\alpha{=}0.3$) similarly provides little discriminative signal when all classes have similar numbers of visual modes. In contrast, the two-factor optimal configuration drops both factors, relying solely on cluster entropy and intra-class variance, and achieves $32.87 \pm 0.29\%$---a $+1.12\%$ improvement over unguided and a $+0.88\%$ improvement over the four-factor design. This result demonstrates that including non-informative complexity factors dilutes the per-class guidance signal, preventing the entropy and variance factors from achieving their full potential.

\paragraph{The two-factor optimal generalizes across scales.}
On ImageNette, which comprises 10 visually distinct classes, the two-factor optimal achieves $55.27 \pm 0.68\%$, outperforming the four-factor design ($52.64 \pm 0.28\%$) by $+2.63\%$ and unguided ($49.99 \pm 0.22\%$) by $+5.28\%$. On ImageNet-100 (100 classes), the two-factor optimal achieves $25.70 \pm 0.07\%$, outperforming the four-factor design ($25.29 \pm 0.27\%$) by $+0.41\%$ with $3.9\times$ lower variance. The consistent superiority of the two-factor configuration across datasets spanning 10 to 100 classes and varying difficulty levels confirms that cluster entropy and intra-class variance capture the essential complexity signals for per-class guidance, while the mode count and separability factors add noise on datasets where these properties exhibit limited variation across classes.

\paragraph{Practical guidance.}
These findings yield a clear practical recommendation: the two-factor configuration $(\alpha{=}0, \beta{=}0.5, \gamma{=}0.5, \delta{=}0)$ should be the default choice, as it provides consistent improvements across all tested datasets without the risk of marginal gains observed with the four-factor design on fine-grained data. The four-factor design may still be preferred when the dataset is known to exhibit high variation in mode count and separability across classes (see the semantic subset analysis in Section~\ref{sec:semantic-subset}), but the two-factor design's robustness and lower variance make it the safer default for unknown datasets. Both approaches remain training-free and require only a one-time complexity analysis of the real data, preserving the practical appeal of the \ags{} framework.
